\lecture{11}{2026-10-22}{Picking Curves}

\begin{example}[P-256, NIST curve]
    The NIST P-256 curve (also known as secp256r1) is defined by the short Weierstrass equation
    \[
        y^2 = x^3 - 3x + b \quad \text{over } \mathbb{F}_p,
    \]
    where \(p = 2^{256} - 2^{224} + 2^{192} + 2^{96} - 1\) and \(b\) is a specific constant.

    The prime \(p\) is a generalized Mersenne prime: its binary expansion is sparse, allowing modular reduction to be computed efficiently via word shifts and additions on 32-bit registers.
    The coefficient \(a = -3\) is chosen for performance: in Jacobian projective coordinates \([X : Y : Z]\), the tangent slope formula in point doubling involves \(3X^2 + aZ^4\); when \(a = -3\), this factors as \(3(X - Z^2)(X + Z^2)\), saving several modular multiplications.

    In fact, NIST/NSA didn't provide \(b\) directly, but instead provided a seed value \(s\) and defined \(b\) as the first 32 bytes of the SHA-1 hash of \(s\).
    The logic is that, even if NIST/NSA figured out a way to pick a "bad" \(b\) that would allow them to break the discrete logarithm problem on this curve, they wouldn't be able to engineer a specific seed \(s\) that hashes to that vulnerable curve (by preimage resistance of SHA-1), unless such weak curves are dense among all elliptic curves.
\end{example}

\begin{example}[Curve25519]
    Introduced by Bernstein~\cite{Bernstein2006}, Curve25519 is defined in \emph{Montgomery form} by the equation
    \[
        y^2 = x^3 + 486662x^2 + x \quad \text{over } \mathbb{F}_p,
    \]
    where \(p = 2^{255} - 19\).

    The group order is \(|E(\mathbb{F}_p)| = 8q\), where
    \[
        q = 2^{252} + 27742317777372353535851937790883648493
    \]
    is a large prime.
    Because the curve has cofactor \(h = 8\), protocols clear the cofactor (e.g., clamping the private key by setting the lowest 3 bits to 0) to ensure scalar multiplications remain in the large prime-order subgroup.
\end{example}

\section*{Invalid Point Attacks}

Recall the group law on an elliptic curve in short Weierstrass form \(y^2 = x^3 + ax + b\).
Given points \(P_1 = (x_1, y_1)\) and \(P_2 = (x_2, y_2)\), the sum \(P_3 = P_1 + P_2 = (x_3, y_3)\) is computed via
\[
    \lambda =
    \begin{cases}
        \frac{y_2 - y_1}{x_2 - x_1} & \text{if } P_1 \neq P_2, \\[0.5em]
        \frac{3x_1^2 + a}{2y_1} & \text{if } P_1 = P_2,
    \end{cases}
    \qquad
    \begin{aligned}
        x_3 &= \lambda^2 - x_1 - x_2, \\
        y_3 &= \lambda(x_1 - x_3) - y_1.
    \end{aligned}
\]
Notice that the addition and doubling formulas do not use the constant \(b\) of the curve at all.

\subsection*{The Attack}

In a Diffie--Hellman exchange, an honest party computes a scalar multiple \(k P\) of an incoming point \(P = (x, y)\).
If an attacker sends an invalid point \(P = (x, y) \in \mathbb{F}_p \times \mathbb{F}_p\) that does \emph{not} lie on the agreed curve \(E\), an honest party blindly applying the group law will nevertheless compute the coordinates of \(k P\).
Because \(b\) does not appear in the addition formulas, this computation actually takes place on a different elliptic curve:
\[
    E': y^2 = x^3 + ax + b', \quad \text{where } b' = y^2 - x^3 - ax \pmod p.
\]
If the attacker carefully selects \((x, y)\) such that \(E'(\mathbb{F}_p)\) has smooth order (or contains a small prime factor \(r\)), the attacker can choose \(P\) of order \(r\) in \(E'(\mathbb{F}_p)\).
When the honest party returns or uses \(k P\), the attacker can determine \(k \pmod r\) using Pohlig--Hellman.
Repeating this with points on several such curves allows the attacker to reconstruct the entire private key \(k\) via the Chinese Remainder Theorem~\cite{BiehlMeyerMueller2000}.

\subsection*{Point Validation and Patents}

An obvious countermeasure is to perform \emph{point validation}: whenever a point \((x, y)\) is received, verify that
\[
    y^2 \equiv x^3 + ax + b \pmod p.
\]
However, this check is easy for implementers to omit.
More problematically, Certicom (later acquired by BlackBerry) held a patent on elliptic curve public key validation (US Patent 6,563,928), meaning implementers risked patent infringement and royalties simply by validating that an incoming point lies on the curve.

\subsection*{Curve25519: $x$-Only Arithmetic and Twist Security}

To completely eliminate point validation and avoid patent issues, Curve25519 uses Montgomery form and \emph{\(x\)-only arithmetic} via the Montgomery ladder.
Because \(P = (x, y)\) and \(-P = (x, -y)\) share the same \(x\)-coordinate, the \(x\)-coordinate of \(k P\) can be computed using only the \(x\)-coordinate of \(P\) and the curve parameter \(A\), without ever referencing or calculating \(y\).
In X25519 key exchange, parties only transmit the 32-byte \(x\)-coordinate.

Could an attacker still mount an invalid point attack by sending an \(x\)-coordinate that does not correspond to any valid point on the curve?
For any field element \(x \in \mathbb{F}_p\), the value \(x^3 + Ax^2 + x\) is either:
\begin{enumerate}
    \item A quadratic residue in \(\mathbb{F}_p\), in which case \(y^2 = x^3 + Ax^2 + x\) has solutions \(y \in \mathbb{F}_p\), meaning \(x\) is the \(x\)-coordinate of a point on \(E(\mathbb{F}_p)\);
    \item Zero, corresponding to a 2-torsion point on \(E(\mathbb{F}_p)\); or
    \item A quadratic non-residue in \(\mathbb{F}_p\).
\end{enumerate}
In the third case, \(x\) does not correspond to a point in \(E(\mathbb{F}_p)\).
However, fixing any quadratic non-residue \(d \in \mathbb{F}_p^\times\), the quotient \((x^3 + Ax^2 + x)/d\) is a non-zero quadratic residue in \(\mathbb{F}_p\).
Thus, there exists \(y \in \mathbb{F}_p\) satisfying
\[
    E^t: d y^2 = x^3 + Ax^2 + x \quad \text{over } \mathbb{F}_p.
\]
The curve \(E^t\) is called the \emph{quadratic twist} of \(E\).
(Over the quadratic extension \(\mathbb{F}_{p^2}\), substituting \(y' = y \sqrt{d}\) maps \(E^t\) isomorphically to \(E\), but they are not isomorphic over \(\mathbb{F}_p\).)

This means \textbf{every field element \(x \in \mathbb{F}_p\) is the \(x\)-coordinate of a valid point on either the curve \(E(\mathbb{F}_p)\) or its quadratic twist \(E^t(\mathbb{F}_p)\)}.
Moreover, the algebraic relations of the Montgomery ladder depend exclusively on \(x\) and \(A\), so the exact same ladder formulas compute scalar multiplication whether \(x\) lies on \(E(\mathbb{F}_p)\) or on \(E^t(\mathbb{F}_p)\).

If an attacker sends an \(x\) on the twist, the honest party's computation implicitly takes place in the twist group \(E^t(\mathbb{F}_p)\).
If \(E^t(\mathbb{F}_p)\) had smooth order, the attacker could recover \(k\) using Pohlig--Hellman on the twist.
By Hasse's Theorem, if \(|E(\mathbb{F}_p)| = p + 1 - t\), the quadratic twist has order
\[
    |E^t(\mathbb{F}_p)| = p + 1 + t.
\]
To protect against this attack, the curve must satisfy \emph{twist security}: its quadratic twist must also have a cryptographically secure group order (a small cofactor times a large prime).

Curve25519 was designed specifically to be twist-secure:
\[
    |E(\mathbb{F}_p)| = 8q, \qquad |E^t(\mathbb{F}_p)| = 4q',
\]
where \(q \approx 2^{252}\) and \(q' = 2^{253} - 55484635554744707071703875581767296995\) are both primes.
The parameter \(A = 486662\) was chosen as the smallest integer with \(A \equiv 2 \pmod 4\) such that both \(E\) and \(E^t\) have near-prime order and \((A - 2)/4 = 121665\) is small for computational efficiency.
As a result, any 32-byte input can be accepted safely: whether it lies on the curve or its twist, the discrete logarithm problem remains equally intractable, rendering invalid point attacks impossible without any need for point validation.